% 建议负责人：负责模型改进与实验分析的成员
\section{创新点}

本节将创新点收束为三条主线：分层混合情感分析、原文证据保留机制、证据约束的建议生成。长文本
滑动窗口和独立压力测试不单独作为同等级创新点，而是作为这些设计的实现支撑和评估证据。

\subsection{创新性设计一：分层混合情感分析}

项目没有将“使用 BERT”本身作为创新点，而是针对课程评价的真实问题设计分层混合流程。BERT
负责主要上下文语义判断；规则模块处理课程场景中可明确描述的转折、否定、建议和混合评价；
TF-IDF Logistic Regression 与纯规则模型承担自动回退。不同方法承担各自适合的任务，而不是
简单投票。

该流程还包含长文本滑动窗口推理。系统没有直接截断超过 160 Token 的评价，而是使用带重叠的
窗口，并按每个窗口的有效 Token 数加权聚合概率。窗口数量、Token 数和截断标志会传递到前端，
使教师能够看到超长文本是否触及处理上限，避免后台静默丢弃信息。

12 条固定案例的消融结果见图~\ref{fig:ablation-visualization}。这些案例主要用于验证课堂演示
流程，不等同于独立泛化测试。纯规则取得满分，是因为案例集中包含较多规则能够直接覆盖的典型
表达；因此项目另外使用未参与规则开发的 48 条压力测试评估复杂输入下的鲁棒性。

\subsection{创新性设计二：原文证据保留机制}

多数文本清洗流程会直接修改整个输入，从而使后续证据片段与学生原话不一致。本项目将“用于分类
的标准化文本”和“用于展示与证据的原始文本”分开。缩写展开、重复字母收缩和拼写纠错只进入
情感分类通道；课程维度、关键词、相似评价和证据仍从原文产生。

同时，纠错策略限制在情感词范围内，并要求极小编辑距离与唯一候选，避免将课程名、编程库名称
或合法英文词形错误修改。这种设计在提高噪声文本鲁棒性的同时，保留了分析结果的可追溯性。
教师可以从系统输出回到学生原话，判断系统是否正确理解了“作业任务”“考试安排”或“实验实践”
等具体维度。

\subsection{创新性设计三：证据约束的建议生成}

系统把大语言模型限制在“表达与归纳”环节，而不是让它直接决定情感和课程问题。提示词要求建议
中的问题维度优先来自已识别的课程维度，证据必须引用原始评论、主题证据或相似评价。即使 LLM
不可用，本地模板也会输出相同字段。这种结构降低了生成式模型幻觉对核心分析结论的影响。

从技术贡献看，本项目主要提高了课程评价系统在常见噪声和混合表达下的工程鲁棒性。独立压力测试
包含 48 条样本，每类 6 条，覆盖缩写、拼写错误、否定、讽刺、隐含抱怨、中英混合、长文本和
混合情感。完整混合流程通过 34/48 条样本，Accuracy 为 0.7083，Macro-F1 为 0.7217，均高于
rule-only、TF-IDF-only 和 BERT-only。另一方面，讽刺类别仅通过 1/6，隐含抱怨仅通过 3/6。
这些失败类型被保留在报告中，用于明确后续数据扩充和模型改进方向。
